\section{Fisher Grouping and Moment Proofs}
\label{annexsec:p15-fisher-grouping}

These derivations refer to the explicitly named recording experiments in
\cref{sec:p15-fisher-grouping}. An aperiodic average, a fixed-grid moment
and an asymptotic sampling variance are different objects.

\begin{annexitem}{annex:p15.fg.phase}{thm:p15-fg-phase}{Phase averaging and inverse corrections}
\begin{proof}
For fixed \(x\), as the phase runs once through \([0,2\pi)\), the
midpoint assigned to \(x\) runs once, with constant speed \(h/(2\pi)\),
through \([x-h/2,x+h/2]\), apart from a single endpoint convention.
Thus its displacement from \(x\) has the uniform law on that interval
shifted to zero. For every integrable function \(b\),
\[
 \frac1{2\pi}\int_0^{2\pi}b(Y_\vartheta(x))d\vartheta
 =\frac1h\int_{-h/2}^{h/2}b(x+u)du.
\]
The bound \(|Y_\vartheta|\leq|X|+h/2\), together with
\((r+s)^k\leq2^{k-1}(r^k+s^k)\) for \(k\geq1\), supplies an
integrable bound for the absolute \(k\)-th moment. Fubini therefore
applies to the displayed identity with \(b(y)=y^k\). It gives the joint
law representation \(X+U\), with \(U\) independent uniform error;
independence here concerns the phase-averaged experiment only.

Direct integration gives \(\mathbb EU=\mathbb EU^3=0\),
\(\mathbb EU^2=h^2/12\), and \(\mathbb EU^4=h^4/80\).
Expanding \((X+U)^k\) proves all four forward identities. Solve the
first three triangular equations. In the fourth, substitute
\(M_2=\overline G_2-h^2/12\), obtaining
\[
 M_4=\overline G_4-\frac{h^2}{2}\overline G_2
       +h^4\left(\frac1{24}-\frac1{80}\right)
     =\overline G_4-\frac{h^2}{2}\overline G_2+\frac{7h^4}{240}.
\]
For the uniform counterexample the fixed grid records zero almost surely;
direct integration gives \(\mathbb EX^2=h^2/12\). Hence the negative
fixed-grid corrected value is indeed a failure, not an undefined moment.
\end{proof}
\end{annexitem}

\begin{annexitem}{annex:p15.fg.fourier}{thm:p15-fg-fourier}{Normal Fourier coefficients and convergence}
\begin{proof}
Put \(g(y)=\int_{y-h/2}^{y+h/2}f(x)dx\). By definition
\(G_k(\vartheta)=\sum_jc_j^kg(c_j)\). For a normal density,
the summands and each phase derivative are uniformly summable: on
bounded phase intervals their absolute values are bounded by a polynomial
in \(|j|\) times a constant times \(e^{-d j^2}\), for some \(d>0\).
This follows from differentiating normal densities at the bin endpoints;
each derivative is a polynomial times that density. It justifies the
termwise integrations below and makes \(G_k\) a smooth periodic function.

Its complex Fourier coefficient of index \(s\) is
\[
 \frac1{2\pi}\int_0^{2\pi}G_k(\vartheta)e^{-is\vartheta}d\vartheta
 =\frac1h\int_{\mathbb R}y^kg(y)e^{-2\pi isy/h}dy.
\]
Indeed substitute \(y=h(j+\vartheta/(2\pi))\) in each summand;
the intervals \([hj,h(j+1)]\) tile the real line, and
\(e^{-is\vartheta}=e^{-2\pi isy/h}\).
All absolute moments of \(g\) are finite because its integral is \(h\)
and the uniform convolution has bounded error. Consequently
\[
 \int y^kg(y)e^{-ity}dy=i^k\frac{d^k}{dt^k}
                       \int g(y)e^{-ity}dy
\]
by dominated differentiation, with dominating function \(|y|^kg(y)\).

For completeness the centered Gaussian transform is
\(F(t)=e^{-\sigma^2t^2/2}\). Differentiation under its integral is
dominated by \(|x|f(x)\); integration by parts using
\(f'(x)=-xf(x)/\sigma^2\) gives \(F'(t)=-\sigma^2tF(t)\).
The boundary terms vanish by Gaussian decay, and \(F(0)=1\), so the
unique solution of this elementary linear differential equation is the
stated transform. Interchanging the absolutely integrable convolution
integrals gives
\[
 \int g(y)e^{-ity}dy
 =F(t)\int_{-h/2}^{h/2}e^{-itu}du
 =e^{-\sigma^2t^2/2}\frac{2\sin(ht/2)}t=H(t).
\]
This proves the coefficient formula, including its continuous value at zero.

At nonzero lattice frequencies, every derivative of \(H\) is a finite
sum of Gaussian derivatives times derivatives of \(2\sin(ht/2)/t\).
The former are polynomials times \(e^{-\sigma^2t^2/2}\), and the latter
are finite sums of bounded trigonometric factors times nonpositive powers
of \(|t|\). Thus \(|C_{k,s}|\leq A_k(1+s)^{b_k}e^{-cs^2}\)
for suitable finite constants and integer \(b_k\). This proves absolute
uniform summability of the proposed series.

To verify reconstruction rather than assume it, subtract that uniformly
convergent series from \(G_k\). The difference is continuous periodic
and has all Fourier coefficients zero. Its convolution with
\[
 K_r(t)=\frac1r\left|\sum_{j=0}^{r-1}e^{ijt}\right|^2
\]
is therefore zero: expanding the finite square expresses the convolution
as a finite weighted sum of those coefficients. The kernel is nonnegative,
has integral \(2\pi\), and satisfies
\(K_r(t)\leq1/[r\sin^2(\delta/2)]\) when the circular distance of
\(t\) from zero is at least \(\delta>0\). Uniform continuity controls
the convolution error on the remaining \(\delta\)-neighborhood; this
bound controls its complement. Hence these convolutions converge uniformly
to the difference, which must be zero. The reconstruction is proved.

At \(\omega_s=2\pi s/h\),
\(H'(\omega_s)=e^{-cs^2}h(-1)^s/\omega_s\).
Multiplying by \(i/h\) and taking twice the real part of
\(C_{1,s}e^{is\vartheta}\) gives
\((-1)^{s+1}h e^{-cs^2}\sin(s\vartheta)/(\pi s)\).
Also
\[
 H''(\omega_s)=e^{-cs^2}(-1)^s
       \left(-2\sigma^2h-\frac{2h}{\omega_s^2}\right),
\]
which gives the displayed second-moment cosine coefficient. Taking three
and four derivatives is the exact higher-moment rule. For a truncation at
\(m\), its general error is bounded by
\(2h^{-1}\sum_{s>m}|H^{(k)}(\omega_s)|\); the already proved
Gaussian-polynomial bound makes this a convergent explicit tail, uniformly
in phase. For \(k=1\), use \(s^{-1}\leq(m+1)^{-1}\) and
\(s^2\geq(m+1)^2+(s-m-1)(2m+3)\). The resulting geometric series
is the asserted sharper tail bound.

Finally setting \(h=\sigma\) in the first principal amplitude and
solving \(\sigma/(10\sqrt n)=\sigma e^{-2\pi^2}/\pi\)
gives \(n=\pi^2e^{4\pi^2}/100\). The other three source comparisons
follow by the same squaring step from respective amplitude/standard-
deviation pairs
\(4\sigma^2e^{-2\pi^2},\sqrt2\sigma^2/\sqrt n\),
\(12\pi\sigma^3e^{-2\pi^2},\sqrt{15}\sigma^3/\sqrt n\), and
\(32\pi^2\sigma^4e^{-2\pi^2},\sqrt{96}\sigma^4/\sqrt n\).
These pairs are the source's leading approximations, not the complete
coefficients just proved.
\end{proof}
\end{annexitem}

\begin{annexitem}{annex:p15.fg.moments}{thm:p15-fg-moments}{Moment existence and sampling precision}
\begin{proof}
The two inequalities \(|Y_\vartheta|\leq|X|+h/2\) and
\(|X|\leq|Y_\vartheta|+h/2\) show equivalence of every positive
absolute moment. For \(r\geq1\) use the convex bound already given;
for \(0<r\leq1\), use \((u+v)^r\leq u^r+v^r\).
Let \(Z=Y_\vartheta^k\). When \(\mathbb E|Z|<\infty\), set
\(Z^{(b)}=\max(-b,\min(Z,b))\). Then
\(\mathbb E|Z-Z^{(b)}|\to0\) by dominated convergence.
The average of independent \(Z_i^{(b)}\) has variance at most \(b^2/n\)
and hence converges in probability to \(\mathbb EZ^{(b)}\), by the
elementary variance bound
\(P(|W-\mathbb EW|>t)\leq\operatorname{Var}(W)/t^2\).
The expected absolute difference between the truncated and untruncated
averages is at most \(\mathbb E|Z-Z^{(b)}|\), independently of \(n\).
Its probability of exceeding any \(t>0\) is at most that quantity
divided by \(t\), by integrating the indicator of the exceeding event.
First choose \(b\) large, then \(n\) large; the triangle inequality,
including the difference of expectations, proves convergence in probability.

If \(\mathbb EZ^2<\infty\), independence makes the covariance of
distinct repetitions zero; expanding the square gives
\(\operatorname{Var}(n^{-1}\sum_iZ_i)=\operatorname{Var}(Z)/n\).
For the stated Pareto density, direct integration gives
\(\mathbb EX^r=5/(5-r)\) for \(r<5\), and divergence for
\(r\geq5\). Thus \(\mathbb EX^4=5\) and \(\mathbb EX^8=\infty\).
For any fixed grid the same existence conclusions hold for \(Y\).
Since \(Y^4\geq0\), the squared sample average is at least
\(Y_1^8/n^2\), and its expectation is infinite. Its mean is finite,
so its variance is infinite. This does not affect the truncation proof
of consistency, which required only the finite first absolute moment of
\(Z=Y^4\).
\end{proof}
\end{annexitem}

\begin{annexitem}{annex:p15.fg.deviance}{thm:p15-fg-deviance}{Likelihood, score and remainder}
\begin{proof}
The multinomial probability is
\(N!\prod_jp_j^{n_j}/\prod_jn_j!\). Its log, apart from a
data-only constant, is \(\sum_jn_j\log p_j\). Since
\(m_j=Np_j\), subtracting this from
\(\sum_jn_j\log(n_j/N)\) gives exactly \(D\).
The convention at zero counts follows by continuity of \(u\log u\).
Differentiate the finite sum to obtain
\(\partial_\eta D=-\sum_j(n_j/m_j)\partial_\eta m_j\).
Expanding \(Q=\sum_jn_j^2/m_j-N\), using both sums of counts
equal to \(N\), gives its derivative as stated. For countably many
cells, the same likelihood ratio is obtained from the finitely many
occupied cells when its logarithm exists; none of this differentiates an
uncontrolled infinite series.

Write \(n_j=m_j(1+\delta_j)\). The identity
\(\sum_jm_j\delta_j=0\) gives
\(D=\sum_jm_j h_0(\delta_j)\), where
\(h_0(u)=(1+u)\log(1+u)-u\). It has
\(h_0(0)=h_0'(0)=0\), \(h_0''(0)=1\), and
\(h_0'''(u)=-(1+u)^{-2}\). Taylor's formula with integral
remainder therefore gives, for \(|u|\leq\rho<1\),
\[
 \left|h_0(u)-\frac{u^2}{2}\right|
 \leq\frac{|u|^3}{6(1-\rho)^2}.
\]
Multiply by \(2m_j\), sum, and use
\(\sum_jm_j|\delta_j|^3\leq\rho\sum_jm_j\delta_j^2=\rho Q\).
This proves the bound. Differentiating once more for the third-order
term, or expanding \(h_0'''(0)=-1\), gives the reported cubic term
\(-\sum_jm_j\delta_j^3/3\) in \(2D\). No probability distribution
for \(Q\) was used or obtained.
\end{proof}
\end{annexitem}

\begin{annexitem}{annex:p15.fg.information}{thm:p15-fg-information}{Exact Gaussian cell derivatives and information loss}
\begin{proof}
The endpoint formula for \(p_j\) follows by substituting
\(z=(x-\mu)/\sigma\) in its defining integral. The chain rule,
\(\Phi'=\phi\), \(\partial_\mu z_j=-1/\sigma\) and
\(\partial_\sigma z_j=-z_j/\sigma\), gives both derivatives.
Infinite endpoints are handled by the limits \(\phi(\pm\infty)=0\)
and \(z\phi(z)\to0\).

The normal scores are
\[
 S_\mu(x)=\frac{x-\mu}{\sigma^2},\qquad
 S_\sigma(x)=\frac1\sigma\left(\frac{(x-\mu)^2}{\sigma^2}-1\right).
\]
On each compact parameter set with scale bounded away from zero,
the density derivatives are bounded by an integrable polynomial times
a Gaussian with a fixed slower decay. This permits differentiation under
each fixed-cell integral. Consequently
\(\partial_\eta p_j=\int_{B_j}S_\eta(x)f_{\mu,\sigma}(x)dx\),
and dividing by the strictly positive \(p_j\) proves
\(\partial_\eta\log p_j=\mathbb E[S_\eta\mid B_j]\).
The Gaussian moments, found recursively by integrating
\(z^{r-1}(-\phi'(z))\), satisfy
\(\mathbb EZ^{2r}=(2r-1)\mathbb EZ^{2r-2}\) and odd moments zero.
They give zero expected scores and finite score second moments.

For any square-integrable scalar \(S\), expand
\(S=(S-\mathbb E[S\mid B])+\mathbb E[S\mid B]\).
The cross term has expectation zero on every cell, by the defining
conditional integral. Summing nonnegative within-cell integrals is
licensed by countable additivity, and yields
\[
 \mathbb ES^2=\mathbb E(\mathbb E[S\mid B])^2
                   +\mathbb E\operatorname{Var}(S\mid B).
\]
Applying this to the scores proves the scalar identity. Applying it to
every real linear combination of the two scores proves the positive-
semidefinite matrix statement; no scalar inversion with an ignored
nuisance parameter is involved.

The historical midpoint calculation is finite: multiplication of the
nine counts by \((j+1/2)^2\) gives the terms
\(28.5,189,331.25,294,283.5,181.5,126.75,56.25,72.25\), whose sum
is \(1563\). Their count sum is \(300\). Hence the uncorrected
moment is \(521/100\), and the corrected one is
\(521/100-1/12=769/150\), as stated in the body. These arithmetic
identities do not certify the source's separately interpolated likelihood
root or minimum-\(\chi^2\) result.
\end{proof}
\end{annexitem}

\begin{annexitem}{annex:p15.fg.expansion}{thm:p15-fg-expansion}{Controlled phase-averaged information and coefficient qualification}
\begin{proof}
Fix a physical width \(h\). In the independently phase-randomized
experiment the recorded midpoint has density
\(\sigma^{-1}v_a(z)\), where \(a=h/\sigma\) and
\(z=(y-\mu)/\sigma\). It integrates to one by Fubini.
For a fixed grid the cell probability is \(a v_a(z_j)\);
its location/scale derivatives, at fixed physical endpoints, are the same
logarithmic derivatives of \(\sigma^{-1}v_a(z_j)\). Averaging the
grid sum over phase replaces \(a\sum_j\) by \(\int dz\),
by the substitution used in the Fourier proof. This explains exactly
which aperiodic regime is under consideration.

Write \(q=\partial_z\log v_a\) and
\(t=1+zq+a\partial_a\log v_a\). Chain differentiation gives scores
\(-q/\sigma\) and \(-t/\sigma\). Thus
\[
 J_\mu=\int_{\mathbb R}v_a q^2 dz,
 \qquad J_\sigma=\int_{\mathbb R}v_a t^2 dz.
\]
The cross information is zero because \(v_a\) and \(t\) are even in
\(z\), whereas \(q\) is odd. This assertion is for the averaged
experiment; an arbitrary fixed grid need not share this symmetry.

We first justify the expansion and its integrated remainder. Put
\[
 R(a,z)=\frac{v_a(z)}{\phi(z)}
       =\int_{-1/2}^{1/2}e^{-azs-a^2s^2/2}ds.
\]
It is smooth and even in \(a\). Jensen's inequality, or the elementary
convex tangent inequality for the exponential followed by integration,
gives \(R(a,z)\geq e^{-a^2/24}\). On \(|a|\leq a_0\), every
fixed-order derivative in \(a,z\) of the integrand is bounded by
a polynomial in \(1+|z|\) times a constant times
\(e^{a_0|z|/2}\). The derivative follows by repeated product/chain
rules and \(|s|\leq1/2\). Its integral obeys the same bound.
Derivatives of \(1/R\) are finite sums of products of these derivatives
divided by powers of \(R\), whose lower bound is uniform. Therefore
the derivatives through order ten in \(a\) of
\(v_aq^2\) and \(v_at^2\) are bounded by an integrable function
of the form
\(C(1+|z|)^b e^{-z^2/2+C'|z|}\).
The same argument covers the additional mixed derivatives required by
the scores. Completing the square proves this function integrable.
Dominated differentiation under the information integrals is justified.
Taylor's integral remainder of order nine is bounded by
\(|a|^{10}/10!\) times the supremum of the integrated tenth derivative
on the compact interval. Evenness removes all odd coefficients. This
establishes the stated \(O(a^{10})\) regime without integrating an
uncontrolled pointwise expansion in the Gaussian tails.

Here is the finite coefficient calculation. Let \(H_0=1,H_1=z\) and
\(H_{r+1}=zH_r-rH_{r-1}\). Induction from \(\phi'=-z\phi\)
gives \(\phi^{(r)}=(-1)^rH_r\phi\). Integrating the even Taylor
terms across the bin gives
\[
 R=1+\sum_{j=1}^{4}r_{2j}a^{2j}+O(a^{10}),\qquad
 r_{2j}=\frac{H_{2j}(z)}{2^{2j}(2j)!(2j+1)}.
\]
This coefficient assertion may be read as differentiation at zero;
the preceding dominated argument, not a uniform polynomial tail in \(z\),
licenses its use inside the information integral. The first Hermite
polynomials are
\[
\begin{aligned}
 H_2&=z^2-1,\quad H_4=z^4-6z^2+3,\\
 H_6&=z^6-15z^4+45z^2-15,\\
 H_8&=z^8-28z^6+210z^4-420z^2+105.
\end{aligned}
\]
For \(\log R=\sum_{j=1}^4\ell_{2j}a^{2j}+O(a^{10})\), expansion
of \(\log(1+u)\) as \(u-u^2/2+u^3/3-u^4/4\) gives
\[
\begin{aligned}
 \ell_2&=(z^2-1)/24,\\
 \ell_4&=-(z^4+4z^2-2)/2880,\\
 \ell_6&=(z^6+6z^4+3z^2-1)/181440,\\
 \ell_8&=-z^8/9676800-z^6/1209600-z^4/806400
                  +z^2/3628800-1/14515200.
\end{aligned}
\]
These expressions follow respectively from
\(r_2\), \(r_4-r_2^2/2\),
\(r_6-r_2r_4+r_2^3/3\), and
\(r_8-r_2r_6-r_4^2/2+r_2^2r_4-r_2^4/4\).
This lists all the finite algebra, rather than presupposing Sheppard's rule.

For clarity, set \(r_0=1\),
\(q_0=-z\), \(q_{2j}=\ell_{2j}'\),
\(t_0=1-z^2\), and
\(t_{2j}=z\ell_{2j}'+2j\ell_{2j}\).
The coefficient of \(a^{2d}\) in each information integral is the
finite sum
\[
 \sum_{u+v+w=d}\int\phi(z)r_{2u}(z)b_{2v}(z)b_{2w}(z)dz,
 \quad b=q\text{ or }t.
\]
Every index is between zero and four. Odd monomials integrate to zero;
the coefficient of each even monomial \(z^{2e}\) is multiplied by
\((2e)!/(2^e e!)\), from the integration-by-parts recurrence already
proved. Substituting the explicitly listed polynomials gives
\[
\begin{array}{c|rrrrr}
 2d&0&2&4&6&8\\\hline
 J_\mu&1&-1/12&1/144&-1/1728&31/518400\\
 J_\sigma&2&-1/3&1/20&-1/135&83/64800.
\end{array}
\]
The finite-sum formula, polynomials and monomial integrals specify each
entry without a numerical approximation. Multiplying the asserted inverse
series by these rows verifies its coefficients through order eight; the
constant terms are positive, so the controlled remainder and continuity
also justify inversion on a sufficiently small interval.

One can see the precise sixth-order location discrepancy without the
eighth-order table. The preceding polynomials give
\[
 \partial_z^2\log v_a=-1+\frac{a^2}{12}
  -\frac{a^4(3z^2+2)}{720}
  +\frac{a^6(5z^4+12z^2+1)}{30240}+O(a^8).
\]
In this same aperiodic experiment,
\(\mathbb Ez^2=1+a^2/12\) and
\(\mathbb Ez^4=3+a^2/2+a^4/80\), by the phase-average proof.
The sixth-order coefficient of its expectation is therefore
\[
 -\frac1{2880}+\frac{5\cdot3+12\cdot1+1}{30240}
 =-\frac1{2880}+\frac1{1080}=\frac1{1728}.
\]
Integration by parts gives
\(J_\mu=-\int v_a\partial_z^2\log v_a\,dz\), with vanishing
Gaussian-tail boundary terms. This matches the information row, whereas
the printed \(1/4320\) would yield the different inverse-series term
reported by Fisher. The mismatch is within the explicitly matched
aperiodic assumptions, not a comparison of an averaged and unaveraged grid.

Finally \(z\) in the phase-randomized experiment is a standard normal
plus independent uniform error of width \(a\). Its variance is
\(1+a^2/12\), and its fourth moment is \(3+a^2/2+a^4/80\).
Subtracting the square of the second moment gives
\(2(1+a^2/6+a^4/360)\). Independence of repetitions divides these
variances by \(n\). Restoring powers of \(\sigma\) proves the exact
moment-variance statements and the algebraic benchmark comparisons in
the body. Nothing in this calculation proves a sampling limit theorem
for a grouped MLE or a fixed-phase information-remainder bound.
\end{proof}
\end{annexitem}
